% Imagen de PostgreSQL
\subsection{Imagen de PostgreSQL}\label{sec:5_1_carga_imagen}

\subsubsection{Descarga de la imagen oficial postgres:16-alpine}
La base de datos no necesita una imagen propia: se usa la imagen oficial de PostgreSQL 16 publicada en Docker Hub, en su variante \texttt{alpine}, que es más pequeña porque está construida sobre Alpine Linux. La imagen ya se descargó al preparar la VPS (sección \ref{sec:2_2_instalacion_podman}); si no, el script la descarga automáticamente la primera vez. Se comprueba que esté disponible y su versión exacta:
\begin{lstlisting}
# podman images docker.io/library/postgres
# podman run --rm docker.io/library/postgres:16-alpine postgres --version
\end{lstlisting}
El segundo comando crea un contenedor temporal (\texttt{-{}-rm} lo borra al terminar) solo para imprimir la versión, que debe ser \texttt{postgres (PostgreSQL) 16.x}.
\captura{bd-01-imagen-postgres}{Imagen postgres:16-alpine disponible y versión de PostgreSQL}
